\documentclass[10pt,a4paper]{amsart}
\usepackage[margin=19mm]{geometry}
\usepackage{amsmath,amssymb,mathtools}
\usepackage[scr=rsfs,cal=euler]{mathalfa}
\usepackage{xfrac}
\usepackage[unicode,hidelinks,bookmarks=false]{hyperref}
\newcommand{\ie}{i.e.\ }
\newcommand{\cf}{cf.\ }
\newcommand{\resp}{respectively\ }
\newcommand{\Subsection}{\subsection}
\providecommand{\nobreakdash}{\nobreak\hbox{-}}
\setlength{\emergencystretch}{2em}
\multlinegap0pt
\usepackage{amssymb}
\newtheorem{theorem}{Theorem}
\title{Local Reflexivity and Duality for Subspace Approximation Schemes}
\author{Daniel Akech Thiong}
\date{Source: arXiv:2609.10460v1 (CC BY 4.0)}
\begin{document}
\maketitle

\noindent\emph{Source and licence.} Local Reflexivity and Duality for Subspace Approximation Schemes. Version arXiv:2609.10460v1 (CC BY 4.0), distributed by arXiv under the Creative Commons Attribution 4.0 International licence, http://creativecommons.org/licenses/by/4.0/, as recorded in the arXiv licence metadata for this exact version and shown on the abstract page https://arxiv.org/abs/2609.10460v1. Full licence notice: \texttt{licenses/arxiv-2609.10460-CC-BY-4.0.txt}.\par\smallskip

\noindent\textit{Editorial scope.} Counted unit: the article's theorem thm:duality (source0.tex lines 119--124). The article's complete original proof of it is delivered with it (source0.tex lines 126--146, one \texttt{proof} environment of 21 lines). No mathematics is written, repaired or re-derived: the statement and the proof are the article's own lines, in the article's own notation, and no retained line is substituted or re-flowed.  No further source-local context is needed: every cross-reference of every retained line resolves inside the two delivered spans.  Nothing was omitted from any retained span, so the package carries no omission record.  No retained line contains a citation, so the wrapper carries no bibliography and no bibliography entry.  No retained line contains a figure environment, an image-inclusion command or drawing code, so no image is delivered and the package declares no asset dependency.  Editorial formatting only, in addition: an AMS \texttt{amsart} preamble that restates the article's own declarations and macros used by the retained lines, namely the environments \texttt{theorem} on separate counters, together with the standard packages those lines need.  The retained environments print in this document as Theorem 1 in order of appearance, whereas the article prints the same items with the numbering of its own full document.  The complete native source of the article as distributed in the arXiv e-print for this exact version is bundled unchanged as \texttt{sources/209917/source0.tex}.
\begin{theorem}[Unconditional Duality for Finite-Dimensional Schemes] \label{thm:duality}
Let $X$ be a Banach space equipped with a generalized subspace approximation scheme $Q = (Q_n)_{n=0}^\infty$ where each $Q_n$ is finite-dimensional. For any bounded linear operator $T \in \mathcal{L}(X)$, the approximation numbers satisfy unconditionally:
\begin{equation}
a_n(T, Q) = a_n(T^{**}, Q^{\perp\perp}).
\end{equation}
\end{theorem}

\begin{proof}
The forward inequality $a_n(T^{**}, Q^{\perp\perp}) \le a_n(T, Q)$ is a standard consequence of the canonical embedding. If $A \in \mathcal{L}(X)$ satisfies $A(X) \subseteq Q_n$, then its bitranspose satisfies $A^{**}(X^{**}) \subseteq Q_n^{\perp\perp}$. Since $\|T^{**} - A^{**}\| = \|T - A\|$, any valid approximant for $T$ induces a valid approximant for $T^{**}$, yielding the inequality.

We now establish the reverse inequality, $a_n(T, Q) \le a_n(T^{**}, Q^{\perp\perp})$, without utilizing local reflexivity. Let $\epsilon > 0$. By definition, there exists an operator $B \in \mathcal{L}(X^{**})$ such that $B(X^{**}) \subseteq Q_n^{\perp\perp}$ and $\|T^{**} - B\| < a_n(T^{**}, Q^{\perp\perp}) + \epsilon$.

Because $Q_n$ is finite-dimensional, it is reflexive, and thus $Q_n^{\perp\perp} = J(Q_n)$ where $J: X \to X^{**}$ is the canonical isometry. Consequently, the range of the bidual approximant is strictly constrained:
\begin{equation}
B(X^{**}) \subseteq J(Q_n) \subset J(X).
\end{equation}
This allows us to define an operator $A: X \to X$ by strict isometric restriction: $A = J^{-1} B J$. The operator $A$ is well-defined, linear, and bounded. Furthermore, its range satisfies $A(X) = J^{-1}(B(J(X))) \subseteq J^{-1}(J(Q_n)) = Q_n$. Therefore, $A$ is a valid approximant for $T$ with respect to the scheme $Q$.

We compute the norm of the approximation error. For any $x \in X$ with $\|x\| \le 1$:
\begin{equation}
\|Tx - Ax\|_X = \|J(Tx - Ax)\|_{X^{**}} = \|T^{**}Jx - BJx\|_{X^{**}} \le \|T^{**} - B\|_{\mathcal{L}(X^{**})} \|Jx\|_{X^{**}}.
\end{equation}
Taking the supremum over the unit ball yields $\|T - A\|_{\mathcal{L}(X)} \le \|T^{**} - B\|_{\mathcal{L}(X^{**})}$. Thus, we have found an operator $A \in \mathcal{L}(X)$ such that $A(X) \subseteq Q_n$ and:
\begin{equation}
\|T - A\| \le \|T^{**} - B\| < a_n(T^{**}, Q^{\perp\perp}) + \epsilon.
\end{equation}
Since $\epsilon > 0$ is arbitrary, $a_n(T, Q) \le a_n(T^{**}, Q^{\perp\perp})$, concluding the proof.
\end{proof}

\end{document}
